\section{Recommendations}
\label{sec:discussion}

\subsection{Controls both communities should report}

Section~\ref{sec:results2} reports three results that read the other way
without the arm beside them. None of the three arms is expensive. None is
currently standard. We propose all three as reporting requirements.

\textbf{Report a no-retrieval arm.} Any defense that modifies how a model uses
retrieved context must report what the same model scores with no retrieved
context at all. Without it, a defense that has stopped reading the document is
indistinguishable from a defense that has learned to resist a poisoned one.
This is one extra arm and it costs one extra pass over the data.

\textbf{Report the abstention rate, not only accuracy.} A method that converts
confident wrong answers into refusals has done something valuable and will
score as inert on accuracy. In an agent setting, where a wrong answer becomes a
wrong action, that difference is the whole point. Reporting where the
probability mass went, rather than only whether it landed on the gold answer,
costs nothing beyond classifying the generations already produced.

\textbf{Report at least two base models.} Holding defense and attack fixed and
changing only the model changed an effect's magnitude, its existence, and once
its sign. A single-model result in this area should be read as a
model-specific result until shown otherwise.

\subsection{Conventions for measurable transfer}

\textbf{State the regime the evaluation actually covers.} Our defense
census codes each method by what its own paper evaluated, and
\PctOneRegime{}\% cover one regime. Many of those papers describe their
contribution in terms broad enough to imply the other. One required sentence
would close this. Name the regime the experiments cover. A reader of the
paper, and the maintainer of a system, then sees the gap without reading the
appendix.

\textbf{Gate claims on a detectability threshold.} Reporting a difference
without the smallest difference the experiment could have detected makes "no
effect" and "we could not tell" look identical. We computed this for every run
and several of our own results moved from the first category to the second.

\textbf{Say whether the attack was adapted to the defense.} A defense that
holds against an attack built for a different target has been shown to resist
that attack and nothing more. The asymmetry matters and papers rarely state it.
Non-adaptive evaluation is conservative evidence for a failure. It is no
evidence at all for a guarantee. Every published evaluation of a
context-reliance defense in our corpus is non-adaptive. We built an adaptive
payload against the one method that had resisted, and
\PctAcdRemovedGradient{}\% of its protection
went (Section~\ref{sec:adaptive}). One sentence naming whether the
attacker knew the defense would let a reader tell which kind of claim they are
reading.

\textbf{Compose across the boundary, and measure the composition.}
Section~\ref{sec:parallel} shows the two communities defending different
stages of one pipeline, and Section~\ref{sec:results} shows the reasoning
stage failing on its own once an adversary writes the input. The obvious
next move is to put an ingestion-stage filter in front of a
reasoning-stage method and see whether the pair holds where each half does
not. We do not know that it works. Nothing in this paper tests a composed
defense, and a filter that passes a fluent falsehood leaves the reasoning
stage exactly where our results find it. That is the point: the experiment
is cheap, both halves are published with released code, and no paper in
\NCorpus{} has run it.

\subsection{Artifacts}

The reconstruction standard made artifact quality measurable as a side effect.
Among the best-cited untested defenses in the cell we studied, two name
empty repositories. Two more ship code without the trained weights it needs.
One runs. Two further execution-stage defenses do not
run reliably on any freely-available open-weight model.

We note this without blame. Releasing a working artifact is harder than it
looks and is rarely rewarded. But the coverage gap this paper measures is
partly an artifact gap. A defense nobody else can run is a defense nobody else
will test. \PctDefUnmatched{}\% of the defenses in our census have no
confirmed test against any named mechanism.

\subsection{Implications for deployment}

Three practical points follow from the results.

Consider a system that uses a knowledge-conflict method to improve answer
quality, and that also retrieves from a corpus an outsider can write to. The
method then operates outside its validated regime. Our measurements say it
is more likely to increase damage than to reduce it.

A defense that holds against an optimized adversarial passage may fail against
a fluent sentence that simply asserts a plausible falsehood. In our
measurements no defense resisted both. The fluent attack is cheaper to produce. It needs no
gradient access and no surrogate retriever, and it does roughly twice the
damage to an undefended model at every dose. Evaluations built only around
optimized payloads overstate protection. Do not call a defense tested on one
payload class robust to context poisoning.

Finally, the channel matters more than the method name. Supply-chain,
cross-modal and tool-metadata have fewer than one named defense per named
mechanism. If an agent accepts third-party components or machine-readable tool
descriptions, almost nothing published defends that route.
